\subsubsection{K-Nearest Neighbors (k-NN)}
\label{subsubsec:eda_knn}
El algoritmo de los K-Vecinos más Cercanos (k-NN) es un método de aprendizaje supervisado basado en instancias. A diferencia de otros modelos, k-NN no construye un modelo explícito durante el entrenamiento, sino que almacena todas las instancias de entrenamiento para realizar la clasificación en el momento de la consulta.



\textbf{Arquitectura y funcionamiento} \\
El principio de k-NN es la proximidad: asume que las muestras de audio con características acústicas similares pertenecerán a la misma categoría.
\begin{enumerate}
    \item \textbf{Cálculo de distancia:} Cuando se presenta una nueva muestra de audio, el algoritmo calcula la distancia (generalmente la distancia euclidiana) entre esa muestra y todas las demás del conjunto de entrenamiento.
    \item \textbf{Selección de vecinos:} Se seleccionan los $k$ ejemplares más cercanos (donde $k$ es un número entero definido por el usuario).
    \item \textbf{Votación mayoritaria:} La muestra se asigna a la clase que sea más frecuente entre sus $k$ vecinos.
\end{enumerate}

\textbf{El papel del parámetro k} \\
La elección de $k$ es crítica: un valor muy pequeño (ej. $k=1$) hace que el modelo sea sensible al ruido en el audio, mientras que un valor muy grande puede suavizar demasiado las fronteras de decisión.

\textbf{Ventajas específicas para detección de deepfakes}
\begin{itemize}
    \item \textbf{Sin suposiciones sobre los datos:} No asume que las características del audio sigan una distribución lineal, lo que le permite capturar patrones complejos en la distribución de artefactos sintéticos.
    \item \textbf{Eficacia con características bien definidas:} Si la extracción de características logra separar bien las clases, k-NN es extremadamente preciso.
    \item \textbf{Aprendizaje incremental:} Es fácil añadir nuevos ejemplos de ataques de audio sintético al sistema sin necesidad de reentrenar todo el modelo.
\end{itemize}